% !TEX program = pdflatex
% Draft v2. Theorem statements are final (two adversarial verification rounds); proofs are sketches,
% each annotated with its source in the project records. Complete proofs live in the campaign archive
% and are referenced per theorem. Computational verification is summarised in Appendix A.
\documentclass[a4paper,11pt]{amsart}

\usepackage{amssymb,amsmath,amsthm}
\usepackage[margin=1.1in]{geometry}
\usepackage[colorlinks=true,linkcolor=blue,citecolor=blue]{hyperref}

\newtheorem{theorem}{Theorem}[section]
\newtheorem{lemma}[theorem]{Lemma}
\newtheorem{corollary}[theorem]{Corollary}
\newtheorem{proposition}[theorem]{Proposition}
\newtheorem{conjecture}[theorem]{Conjecture}
\newtheorem{problem}[theorem]{Problem}
\theoremstyle{definition}
\newtheorem{definition}[theorem]{Definition}
\newtheorem{example}[theorem]{Example}
\theoremstyle{remark}
\newtheorem{remark}[theorem]{Remark}

\newcommand{\ce}{\operatorname{ce}}
\newcommand{\fre}{\operatorname{f_{re}}}
\newcommand{\odd}{\operatorname{odd}}
% pointer to the complete proof in the project records
\newcommand{\proofsrc}[1]{\par\smallskip\noindent\begin{sloppypar}\small\emph{Proof status: #1.}\end{sloppypar}\par\smallskip}

\title[The cycle--edge decomposition number of a graph]{Exact values and structural properties of the cycle--edge decomposition number $\ce(G)$}

\author{The Front184 Collaboration}
\thanks{Draft v2, \today. Every statement below has passed two independent adversarial review rounds, and
all computational claims are summarised in Appendix~A. Proofs given here are sketches; the complete proofs
are maintained in the project records and are referenced per theorem.}

\begin{document}
\begin{abstract}
A \emph{decomposition} of a graph $G$ is a partition of $E(G)$ into edge-disjoint cycles and single edges,
and $\ce(G)$ denotes the minimum possible number of parts. Erd\H{o}s and Gallai proved in 1959 that
$O(n\log n)$ parts always suffice and conjectured that $O(n)$ do; the conjecture is still open, the best
bound being the $O(n\log^* n)$ of Buci\'c--Montgomery, while the sharper suggestion that $n-1$ parts might
suffice fails already on $K_{3,n-3}$. This paper maps where the $n-1$ budget holds exactly and where it
breaks. We determine $\ce(K_n)$ for every $n$; we prove that $\ce$ is exactly additive over blocks, so that
cut-vertices never help, and derive a transfer principle that reduces every linear bound to its
$2$-connected case. For bounded degree, the $n-1$ bound fails at maximum degree six---the graph formed by a
triangle and a $K_{3,4}$ on the same three vertices has $n=7$ and $\ce=7$---while for Eulerian graphs of
maximum degree at most six we prove $\ce(G)\le n-1$ without exception. The latter resolves the six-regular
case of a problem of Akbari--Aloni--Beikmohammadi--Clow, with a strengthened bound: every $6$-regular graph
on $n$ vertices with $n\equiv 0\pmod 3$ satisfies $\ce(G)\le n-2$. The proof rests on a rigidity lemma,
which says that a $4$-regular graph whose $2$-factors are all unions of triangles must be the line graph
$L(B)$ of a cubic bipartite graph, together with an explicit factorisation of such line graphs into two
triangle-free $2$-factors (an explicit form of an observation of Kouider--Sabidussi). We also prove the
cubic case of the same authors' odd-regular conjecture for the relaxed invariant $\fre$, reduce the general
$\Delta\le 6$ question to an explicitly identified $T$-join danger zone, and establish a sharp degeneracy
threshold: $\ce\le n-1$ holds for every $2$-degenerate graph and fails at every degeneracy $d\ge 3$.
\end{abstract}

\maketitle

\tableofcontents

\section{Introduction}\label{sec:intro}

\subsection{The problem}

Let $G$ be a finite simple graph on $n$ vertices and $m$ edges. A \emph{part} is either a single edge or
the edge set of a simple cycle, and a \emph{decomposition} of $G$ is a partition of $E(G)$ into parts. The
minimum number of parts is $\ce(G)$, the cycle--edge decomposition number of $G$. The definition is naive;
the function is not.

Erd\H{o}s and Gallai \cite{EG59} proved that every $n$-vertex graph admits a decomposition into
$O(n\log n)$ parts, and conjectured that $O(n)$ always suffice; the conjecture in this form is recorded in
\cite{EGP66} and is listed as Problem \#184 in Erd\H{o}s' problem collection \cite{EP184}. Erd\H{o}s also
suggested, in \cite{Er71}, that $n-1$ parts might always be enough. Read without edge-disjointness, that
suggestion is a theorem: Pyber \cite{Py85} showed that $n-1$ cycles and edges always \emph{cover} the edges
of a graph. Read with edge-disjointness, it is false, and it fails early: the complete bipartite graph
$K_{3,n-3}$ already forces $(4/3)(n-3)$ parts, and the construction of Erd\H{o}s \cite{E83}, which uses the
complete bipartite graphs $K_{2k+1,\,n-2k-1}$ and generalises an example of Gallai, forces
$(3/2-o(1))n$ parts; both computations are carried out in \cite{BM22}. And yet no graph is known to need
more than $(3/2-o(1))n$ parts. The linear story is wide open. After
Conlon--Fox--Sudakov \cite{CFS14} improved the general upper bound to $O(n\log\log n)$, Buci\'c and
Montgomery \cite{BM22} obtained the current best, $O(n\log^* n)$.

Why should one study the precise question $n-1$, given that it is false? Because the counterexamples are
highly structured, and knowing exactly which families meet the budget and which break it says something
about where the true constant of the linear conjecture comes from. That is the point of view of this paper:
we treat $n-1$ as a budget and audit it, family by family.

\subsection{This paper}

The audit proceeds on four fronts. First, structure: cut-vertices never help. We prove that $\ce$ is
exactly additive over blocks,
$\ce(G)=\sum_{B}\ce(B)$ (Theorem~\ref{thm:blockadd}), so the entire difficulty of the problem lives on
$2$-connected blocks, and every linear bound can be transferred between general graphs and their blocks
(Corollary~\ref{cor:transfer}). At the two ends of the density range the budget is met with zero slack, for
opposite reasons: trees have no cycles to merge, while even-order complete graphs are pinned by parity
--- every decomposition of $K_{2m}$ needs at least $m$ single edges --- and the two accounts close at
exactly $n-1$ (Theorem~\ref{thm:complete}).

Second, bounded degree. The $n-1$ bound holds for $\Delta\le4$ \cite{AABC25}, whose authors ask what
happens at $\Delta=5$ and $\Delta=6$. At $\Delta=6$ the budget breaks: the graph $K_3\cup K_{3,4}$ --- a
$K_{3,4}$ and a triangle on the same three vertices --- has $n=7$, $\Delta=6$, and
$\ce=7=n$ (Theorem~\ref{thm:counterex}). The interesting news is on the Eulerian side. For Eulerian graphs
of maximum degree $6$ we close the problem completely at $n-1$, and this resolves the six-regular case of
Problem 6.1 of \cite{AABC25} with a strengthened bound
\begin{center}
every $6$-regular simple graph with $n\equiv 0\pmod 3$ satisfies $\ce(G)\le n-2$,
\end{center}
together with the verbatim proof of Conjecture 6.2 of \cite{AABC25} (every $6$-regular graph decomposes
into three $2$-factors, one of which has a component of order at least $4$). The obstacle that had held up
the six-regular case is genuinely structural. If all $2$-factors of a $4$-regular graph are unions of
triangles, then the graph is forced to be the line graph $L(B)$ of a cubic bipartite graph
(Lemma~\ref{lem:fourreg}); on such line graphs, two complementary $2$-factors $\Phi(M)$ and $\Phi(-M)$,
built from an orientation of a perfect matching of $B$, partition the edge set into two triangle-free
$2$-factors (Proposition~\ref{prop:antiphi}) --- a contradiction. That last construction is an explicit,
self-contained special case of an observation of Kouider--Sabidussi \cite{KS95}, and we claim no priority
for the existence statement itself; the reduction, the explicit partition, and the resulting bounds are,
to the best of our knowledge, new. The general $\Delta\le6$ question is then reduced to a $T$-join bound
with an honestly identified danger zone (Theorem~\ref{thm:d6tjoin}, Remark~\ref{rem:danger}).

Third, odd degrees on the relaxed invariant. The paper \cite{AABC25} also studies $\fre(G)$, the minimum
number of blocks in a partition of $E(G)$ into $2$-regular subgraphs and single edges, and conjectures
(Conjecture~6.4 there) that every $(2k+1)$-regular graph has $\fre\le n-1$. The case of graphs with a
perfect matching is easy and was already observed in \cite{AABC25}; the open core is the graphs without
one. We prove the cubic case unconditionally: every cubic graph satisfies $\fre(G)\le n-1$
(Theorem~\ref{thm:cubicpmfree}). The argument is a reduction chain --- an exact formula
$\fre=p+1$ on cubic graphs in terms of the minimum $T$-join, a decomposition of $p$ over the block tree,
and a lower bound on even subgraphs of bridgeless subcubic blocks --- whose last step is quoted from the
published Theorem 1.3 of Choi--Kim--Kostochka--Park--West \cite{CKKPW19}; our contribution is the chain
itself and an exhaustive computational verification.

Fourth, sparsity. The degeneracy hierarchy collapses cleanly: $\ce\le n-1$ holds for every
$2$-degenerate graph --- equivalently every $K_4$-minor-free graph --- and fails, by arbitrarily large
margins, at every degeneracy $d\ge3$ (Theorems~\ref{thm:deg2} and \ref{thm:degex}). The extremal family is
the complete bipartite graph $K_{3,q}$, for which we determine $\ce$ exactly for all $q$
(Theorem~\ref{thm:k3q}), so the asymptotic constant of the $3$-degenerate class is pinned to $[4/3,\,5/3]$.

Finally, we dissect the Buci\'c--Montgomery architecture from the inside (Section~\ref{sec:bm}): the
polylogarithmic remainder map is locked to depth $\log^* n$ by a superconstant fixed point; the apparently
better remainder $2^{(\log d)^\theta}$ in fact yields depth $\Theta(\log\log\log n)$, which is
\emph{worse}; and no amount of edge deletion can create the expanders the method needs. These are
architecture-relative statements --- they locate what any improvement must change, not unconditional lower
bounds against all algorithms.

\subsection{Relation to known results}

Pyber's covering theorem \cite{Py85} gives $n-1$ parts when edges may be reused; our block calculus shows
that for the decomposition version nothing is ever gained or lost at a cut-vertex, so the whole difficulty
is $2$-connected. The planar case of the decomposition problem does not appear to have been studied before;
a literature search in September 2026 found no prior work, though we have not re-verified this against the
full literature. The toroidal and positive-genus bounds of Section~\ref{sec:planar} appear to be new, with
the same caveat, as does the $K_4$-minor-free statement of Theorem~\ref{thm:deg2} (flagged as an
unverified-novelty item in our own records). For Theorem~\ref{thm:kst} no novelty is claimed at all: the
lower bounds are known, the computation is the same as in \cite{BM22,E83}, and we restate it only because
our extremal families are built on it. The six-regular factorisation of line graphs into triangle-free
$2$-factors (Proposition~\ref{prop:antiphi}) is an explicit special case of Kouider--Sabidussi
\cite{KS95}; the related algorithmic result of Bertram--Hor\'ak \cite{BH97} decides whether a
$4$-regular graph decomposes into two triangle-free $2$-factors but contains no counting consequence for
$\ce$. Where our extremal mechanisms coincide with the known $(3/2-o(1))n$ construction --- parity-forced
single edges at the dense side of an odd complete bipartite kernel --- we say so in place.

\section{Preliminaries}\label{sec:prelim}

All graphs are finite and simple. For a graph $G$ we write $n=|V(G)|$, $m=|E(G)|$, $\Delta=\Delta(G)$, and
$c(G)$ for the number of connected components. A \emph{part} is a single edge or the edge set of a simple
cycle of length at least $3$.

\begin{definition}\label{def:ce}
A \emph{decomposition} of $G$ is a partition of $E(G)$ into parts. $\ce(G)$ is the minimum number of parts
in a decomposition, with $\ce(G)=0$ when $E(G)=\emptyset$. Let $f(n)=\max\{\ce(G):|V(G)|=n\}$.
\end{definition}

Two identities do most of the elementary work. The first converts decompositions into cycle packing; the
second converts parity into forced single edges.

\begin{lemma}[Budget/savings identity]\label{lem:savings}
For any decomposition with $k$ cycles of lengths $\ell_1,\dots,\ell_k$ and $s$ single edges,
$k+s=m-\sum_i(\ell_i-1)$. Consequently
\[
\ce(G)\;=\;m-\max_{\mathcal{C}}\ \textstyle\sum_{C\in\mathcal{C}}(|C|-1),
\]
where the maximum is over edge-disjoint cycle families $\mathcal{C}$ (all remaining edges are taken as
single-edge parts).
\end{lemma}
\proofsrc{Complete proof in campaign records, \texttt{campaign/lane\_03\_constant.md} \S0 and
\texttt{campaign/lane\_04\_delta5\_savings.md} \S1; elementary.}

\begin{definition}\label{def:parity}
Let $T=T(G)$ be the set of odd-degree vertices. A subgraph $F\subseteq G$ is a \emph{parity subgraph} if
$\deg_F(v)\equiv\deg_G(v)\pmod 2$ for all $v$; equivalently $F$ is a $T$-join in the sense of
Edmonds--Johnson \cite{EJ73}. Let $\tau(G)$ be the minimum number of edges of a parity subgraph, and let
$\odd(G)=|T|$.
\end{definition}

\begin{lemma}[Parity bounds]\label{lem:parity}
In any decomposition the single edges form a parity subgraph; hence $2s\ge\odd(G)$ and $s\ge\tau(G)$ for
the number $s$ of single edges, and $\ce(G)\ge\odd(G)/2$.
\end{lemma}
\proofsrc{Complete proof in campaign records, \texttt{campaign/lane\_09\_blocks.md} Lemma F; elementary
double counting.}

\begin{lemma}[T-join upper bound]\label{lem:tjoin}
For every connected graph $G$,
$\ce(G)\le \tau(G)+\lfloor (m-\tau(G))/3\rfloor\le\lfloor (m+2\tau(G))/3\rfloor$.
Consequently, if $m+2\tau(G)\le 3n-3$ then $\ce(G)\le n-1$; and $\tau(G)\le n-1$ always, via a spanning
tree, so the criterion applies to every family with $m$ bounded by $3n+O(1)$.
\end{lemma}
\begin{proof}[Proof sketch]
Take a minimum parity subgraph $F$ with $|F|=\tau$. Then $G-F$ is Eulerian, hence decomposes into at most
$\lfloor(m-\tau)/3\rfloor$ edge-disjoint cycles; add the $\tau$ edges of $F$ as single-edge parts. The
criterion is integer arithmetic. For the last claim, the unique minimal $T$-join inside a spanning tree has
at most $n-1$ edges.
\end{proof}
\proofsrc{Complete proof in campaign records, \texttt{campaign/lane\_04\_delta5\_savings.md} Thm.~1.2
(``R1'') and \texttt{campaign/lane\_08\_degeneracy.md} Lemmas 3.1--3.2 (independent re-derivation). Note
that the bound is meaningful only for families with $m=O(n)$; used with $m\le dn-\binom{d+1}{2}$
(degenerate graphs, Theorem~\ref{thm:deg3ub}) and with $m\le3n$ (maximum degree six,
Theorem~\ref{thm:d6tjoin}).}

\begin{definition}\label{def:blocks}
A \emph{block} of $G$ is a maximal subgraph with no cut-vertex of itself; every block is a bridge $K_2$ or
$2$-connected, the blocks partition $E(G)$, and any two blocks share at most one vertex. The
\emph{degeneracy} of $G$ is $d=\max_{H\subseteq G}\delta(H)$, and the $d$-core is the unique maximal
subgraph of minimum degree at least $d$. The \emph{orientable genus} $g(G)$ is the minimum genus of an
orientable surface in which $G$ embeds; all embeddings below are cellular. Orientable genus is additive
over blocks \cite{BHKY62}, whereas non-orientable genus is not \cite{SB77}; all genus statements in this
paper are therefore for the orientable invariant.
\end{definition}

\section{Complete graphs and the block calculus}\label{sec:blocks}

\begin{theorem}[Complete graphs]\label{thm:complete}
For all $m\ge2$, $\ce(K_{2m})=2m-1$; for all $m\ge1$, $\ce(K_{2m+1})=m$.
\end{theorem}
\begin{proof}[Proof sketch]
The upper bounds are the classical Walecki constructions: $K_{2k+1}$ decomposes into $k$ edge-disjoint
Hamilton cycles, and $K_{2k}$ into $k-1$ Hamilton cycles plus a perfect matching; see \cite{Lucas1883} and
\cite{BM08}. For the lower bounds, let $n=2k+1$: every part has at most $n$ edges, so
$\ce\ge\lceil m/n\rceil=(n-1)/2$. Let $n=2k$: all $2k$ vertices are odd, so any decomposition contains
$s\ge k$ single edges by Lemma~\ref{lem:parity}; with $k'$ cycles,
$k'\ge(m-s)/(2k)$ and $\ce=s+k'\ge(2k-1)/2+s(2k-1)/(2k)\ge 2k-1=n-1$.
\end{proof}
\proofsrc{Complete proof in campaign records, \texttt{campaign/lane\_09\_blocks.md} Thm.~G (self-contained
lower bounds; designated proof of record). Numerics: $\ce(K_4)=3$, $\ce(K_5)=2$, $\ce(K_6)=5$,
$\ce(K_7)=3$, $\ce(K_8)=7$ by independent exact solvers, that for $K_8$ being a $3.3\times10^7$-step search
log; see Appendix~A.}

The exact values show two opposite ways of sitting exactly at the budget. Trees do it for lack of
material; $K_{2m}$ does it at full density, because parity pins down $n/2$ single edges and the accounting
then closes at exactly $n-1$. Since zero-slack $2$-connected blocks of the second kind exist at every even
order, no uniform per-block discount can exist. The next results make this precise.

\begin{theorem}[Block additivity]\label{thm:blockadd}
For every graph $G$ with block set $\mathcal{B}(G)$,
\[
\ce(G)\;=\;\sum_{B\in\mathcal{B}(G)}\ce(B).
\]
Moreover every connected graph satisfies the counting identity $\sum_{B}(|V(B)|-1)=n-1$, and more generally
$\sum_B(|V(B)|-1)=n-c(G)$.
\end{theorem}
\begin{proof}[Proof sketch]
For ($\le$), concatenate optimal decompositions of the blocks; block edge sets are disjoint. For ($\ge$),
observe that every part is $2$-regular-connected or a single edge, hence lies in a unique block --- a cycle
meeting two blocks would force two blocks to share three vertices. Grouping parts by block therefore
yields decompositions of the blocks, so the number of parts is at least the right-hand side. The counting
identity follows by induction on the block--cutvertex tree.
\end{proof}
\proofsrc{Complete proof in campaign records, \texttt{campaign/lane\_09\_blocks.md} Thm.~A and Lemma 1.6
(which includes a recorded instance where a $300$-instance randomized check corrected the identity's first
draft).}

Block additivity is an equality, not an inequality: no decomposition can ever exploit a cut-vertex to save
a part. Two corollaries quantify this.

\begin{corollary}[Transfer principle]\label{cor:transfer}
For $\alpha\ge1$: every $2$-connected graph $H$ satisfies $\ce(H)\le\alpha(|V(H)|-1)$ \emph{if and only if}
every graph $G$ satisfies $\ce(G)\le\alpha(n-c(G))$. In particular, for $\alpha=1$ the $n-1$ problem is
exactly equivalent to its $2$-connected restriction.
\end{corollary}

\begin{corollary}[Excess additivity]\label{cor:excess}
$\ce(G)-(n-c(G))=\sum_{B}\bigl[\ce(B)-(|V(B)|-1)\bigr]$: the excess over the linear budget is the sum of
per-block excesses, with no cancellation between blocks.
\end{corollary}
\proofsrc{Complete proofs of Corollaries~\ref{cor:transfer}--\ref{cor:excess} in campaign records,
\texttt{campaign/lane\_09\_blocks.md} Cor.~B/C.}

\begin{proposition}[Merge criterion]\label{prop:merge}
(i) In any decomposition, two distinct parts can never be merged into one part. (ii) $k\ge2$ parts can be
merged into one part if and only if the union of their edge sets is the edge set of a single cycle.
Consequently the cut-vertex discount is exactly zero: no restructuring across, or at, a cut-vertex ever
saves a part.
\end{proposition}
\begin{proof}[Proof sketch]
A union of two parts is the edge set of a cycle only if it is connected and $2$-regular. Two cycles sharing
a vertex give degree $4$ there, and disjoint ones are disconnected; a cycle plus a single edge gives degree
$3$ at an endpoint; two single edges are not $2$-regular.
\end{proof}
\proofsrc{Complete proof in campaign records, \texttt{campaign/lane\_09\_blocks.md} Thm.~D (minimal
witness: the bowtie, $\ce=2$).}

\begin{theorem}[Cactus graphs]\label{thm:cactus}
Let $G$ be connected with every block a bridge or a cycle, and let $t$ be the number of blocks. Then
$\ce(G)=t=(n-1)-\sum_{C\ \mathrm{cycle\ block}}(|C|-2)$.
\end{theorem}
\proofsrc{Complete proof in campaign records, \texttt{campaign/lane\_09\_blocks.md} Thm.~E.}

\begin{remark}\label{rem:zeroslack}
The two zero-slack mechanisms push in opposite directions, and together they explain why the $n-1$ problem
resists both sparse and dense intuition. Numerically, $\ce(K_{2m})=2m-1$ is attained at every tested even
order up to $8$, while $\ce(K_{2m+1})=m\approx n/2$: odd complete graphs enjoy a factor-$2$ discount, their
even cousins none at all.
\end{remark}

\section{Planar, toroidal, and positive-genus families}\label{sec:planar}

\begin{theorem}[Outerplanar graphs]\label{thm:outerplanar}
Every $n$-vertex outerplanar graph satisfies $\ce(G)\le n-1$.
\end{theorem}
\begin{proof}[Proof sketch]
Reduce to blocks. A $2$-connected outerplanar graph has a Hamilton outer cycle, and its weak dual is a
tree. Peel a leaf face: its boundary is one chord $xy$ plus an outer arc whose internal vertices all have
degree exactly $2$, so the arc together with $xy$ forms a cycle $C^*$ whose deletion leaves an outerplanar
graph. Induction gives $\ce(G)\le1+\ce(G')\le1+(n-k)-1\le n-1$, where $k$ is the number of peeled internal
vertices. The degenerate one-chord case is handled explicitly in the archived proof.
\end{proof}
\proofsrc{Complete proof in campaign records, \texttt{notes/attack\_planar.md} Thm.~A with Lemmas A1--A3;
the proof underwent an adversarial fix round (degree-$2$ property of arc-internal vertices; degenerate
case).}

\begin{theorem}[Eulerian plane graphs]\label{thm:eulerplanar}
Every $n\ge3$-vertex Eulerian plane graph satisfies $\ce(G)\le n-2$; in fact $E(G)$ decomposes into at most
$n-2$ cycles, with no single edges.
\end{theorem}
\begin{proof}[Proof sketch]
An Eulerian graph decomposes into edge-disjoint cycles, each of length at least $3$, and $m\le 3n-6$.
\end{proof}
\proofsrc{Complete proof in campaign records, \texttt{notes/attack\_planar.md} Thm.~B.}

\begin{theorem}[Eulerian maximal plane graphs]\label{thm:eulermaxplanar}
If $G$ is an Eulerian maximal plane graph, then $E(G)$ decomposes into exactly $n-2$ edge-disjoint facial
triangles. This is an existence statement for a structured decomposition; $n-2$ need not be $\ce(G)$ ---
the octahedron has $\ce=2$.
\end{theorem}
\begin{proof}[Proof sketch]
The dual $G^*$ is $3$-regular and bipartite, since a plane graph is Eulerian exactly when its dual is
bipartite. A bipartition class of $G^*$ has $(2n-4)/2=n-2$ faces, and every edge of $G$ lies in exactly one
of them.
\end{proof}
\proofsrc{Complete proof in campaign records, \texttt{notes/attack\_planar.md} Thm.~C (the ``minimality''
claim of a first draft was withdrawn in review).}

\begin{theorem}[Bridgeless cubic graphs]\label{thm:cubic}
Every $n$-vertex bridgeless cubic graph satisfies $\ce(G)\le 5n/6$, which is $\le n-1$ for $n\ge6$.
Cubic graphs with bridges satisfy $\ce(G)\le n-1$.
\end{theorem}
\begin{proof}[Proof sketch]
By the Petersen theorem \cite{Pe91} a bridgeless cubic graph has a perfect matching $M$; then $G-M$ is a
$2$-factor with at most $n/3$ cycles. Decompose into $M$ as single edges plus those cycles. Planarity is
not used. The bridged case follows from $\Delta\le4$ \cite{AABC25}.
\end{proof}
\proofsrc{Complete proof in campaign records, \texttt{notes/attack\_planar.md} Thm.~D/D$'$ and
\texttt{campaign/lane\_06\_genus.md} L6.4.}

\begin{theorem}[Eulerian toroidal graphs]\label{thm:eulertorus}
Every Eulerian graph embeddable in the torus satisfies $\ce(G)\le n-1$; moreover any triangulation
component satisfies $\ce\le n-2$.
\end{theorem}
\begin{proof}[Proof sketch]
Work per component. If $e\le 3n-1$, the Euler decomposition gives at most $\lfloor e/3\rfloor\le n-1$
cycles. If $e=3n$, the component is a torus triangulation, hence simple, $2$-connected, and of even
minimum degree at least $4$; by Dirac's circumference theorem \cite{Di52} it contains a cycle of length
$\ge\min\{n,2\delta\}\ge7$, and removing it leaves an Eulerian graph decomposable into at most
$\lfloor(3n_i-7)/3\rfloor=n_i-3$ cycles, for a total of at most $n_i-2$.
\end{proof}
\proofsrc{Complete proof in campaign records, \texttt{campaign/lane\_06\_genus.md} L6.3. The $g=0$
specialisation gives only the weaker $n-1$, so this theorem is compatible with, and does not subsume, the
open planar conjecture $\ce\le n-1$.}

\begin{theorem}[Explicit genus bounds]\label{thm:genus}
Let $G$ have orientable genus $g$.
(i) $\ce(G)\le (m+2n-2)/3\le 5n/3+2g-8/3$.
(ii) If moreover $G$ is Eulerian, then $\ce(G)\le\lfloor m/3\rfloor$; in particular $\ce(G)\le n+2g-3$,
and $\ce(G)\le n+2g-4$ if $G$ triangulates the genus-$g$ surface.
\end{theorem}
\begin{proof}[Proof sketch]
(i) Greedily peel edge-disjoint cycles until a forest $R$ remains; peeling $k$ cycles of total length $L$
costs $k+|E(R)|\le L/3+|E(R)|\le m/3+2(n-1)/3$, and $m\le 3n-6+6g$ by Euler's formula. (ii) The Euler
decomposition has at most $\lfloor m/3\rfloor$ cycles; the triangulation case gains one via the Dirac-cycle
argument of Theorem~\ref{thm:eulertorus}.
\end{proof}
\proofsrc{Complete proof in campaign records, \texttt{campaign/lane\_06\_genus.md} L6.1/L6.2. Part (i) is
a folkloristic observation recorded for completeness.}

\begin{corollary}[Eulerian plane triangulations]\label{cor:eulertripy}
For $n\ge8$, every Eulerian plane triangulation satisfies $\ce(G)\le n-4$. No Eulerian plane triangulation
exists for $n=7$: its degree sequence would be $(6^1,4^6)$, giving a universal vertex whose link is a
$6$-cycle with the remaining three edges drawn outside as a non-crossing perfect matching, that is, a
$3$-regular outerplanar graph --- contradicting the fact that every outerplanar graph has two vertices of
degree at most $2$.
\end{corollary}
\proofsrc{Complete proof in campaign records, \texttt{campaign/lane\_06\_genus.md} L6.2$'$ as revised in
the review round; an earlier ``spoke counting'' argument there contained a counting error and was replaced
by the degree-sequence argument. Exhaustive enumeration double-checks $n\le7$, with the octahedron unique
at $n=6$ (Appendix~A).}

\begin{corollary}[Hamiltonian graphs of genus $g$]\label{cor:hamilton}
Let $G$ have orientable genus $g$ and contain a Hamilton cycle. Then
$\ce(G)\le 1+(m-n)/3+2(n-1)/3\le 4n/3+2g-5/3$. In particular, every $5$-connected toroidal graph --- which
is Hamiltonian by Thomas--Yu \cite{TY97} --- satisfies $\ce(G)\le 4n/3+1/3$.
\end{corollary}
\begin{proof}[Proof sketch]
Delete the Hamilton cycle and apply Theorem~\ref{thm:genus}(i) to the remainder; the last inequality uses
$m\le 3n-6+6g$.
\end{proof}
\proofsrc{Complete proof in campaign records, \texttt{campaign/lane\_06\_genus.md} L6.5.}

\begin{corollary}[A necessary growth rate in $g$]\label{cor:genuslb}
Any universal bound of the form ``every genus-$g$ graph has $\ce\le n-1+f(g)$'' requires
$f(g)\ge\lceil(4g-4)/3\rceil$ for $g\ge2$: the graph $K_{3,\,4g+2}$ has genus exactly $g$ by Ringel's
formula \cite{Ri74}, has $n=4g+5$, and satisfies $\ce(K_{3,4g+2})=\lceil(16g+8)/3\rceil$ by
Theorem~\ref{thm:k3q}.
\end{corollary}
\proofsrc{Complete proof in campaign records, \texttt{campaign/lane\_06\_genus.md} L6.6 (assembled from
Theorem~\ref{thm:k3q}, verified for all $q$, plus Ringel's formula, whose genus-$2$ window was checked
exactly).}

\section{Dirac-regular graphs}\label{sec:dirac}

\begin{theorem}[Dirac-regular graphs]\label{thm:dirac}
Let $n_0$ be the constant of the Hamilton decomposition theorem of Csaba--K\"uhn--Lo--Osthus--Treglown
\cite{CKLOT16}, and let $G$ be a $D$-regular graph on $n\ge n_0$ vertices with $D\ge\lfloor n/2\rfloor$.
Then:
\begin{enumerate}
\item[(i)] if $n$ is odd, $\ce(G)=D/2\le (n-1)/2$, with no single edges;
\item[(ii)] if $n$ and $D$ are even, $\ce(G)=D/2\le (n-2)/2$, with no single edges;
\item[(iii)] if $n$ is even and $D$ is odd, $\ce(G)=(D-1)/2+n/2\le n-1$.
\end{enumerate}
In particular every Dirac-regular graph on $n\ge n_0$ vertices satisfies $\ce(G)\le n-1$.
\end{theorem}
\begin{proof}[Proof sketch]
By \cite{CKLOT16}, $G$ decomposes into Hamilton cycles plus at most one perfect matching. Compare edge
counts modulo $n$: $e(G)=Dn/2$ is divisible by $n$ exactly when $D$ is even, which forces the matching to
be empty in cases (i)--(ii) (in case (i) no perfect matching exists at all) and forces it present with
$t=(D-1)/2$ cycles in case (iii). Maximal $D$ in each regime gives the displayed inequalities; (iii) at
$D=n-1$ with $n$ even gives exactly $n-1$.
\end{proof}
\proofsrc{Complete proof in campaign records, \texttt{campaign/lane\_07\_mindegree.md} Thm.~A, checked
against the local copy of the source of \cite{CKLOT16}.}

\begin{remark}[Scope]\label{rem:diracscope}
The restriction $n\ge n_0$ is essential and must not be dropped: $n_0$ is an astronomical constant, and the
small-$n$, $D\in\{5,6\}$ regular range is not covered by Theorem~\ref{thm:dirac}. No contradiction with the
open status of $\Delta\le5$ follows --- but neither is that case subsumed. The gap between
Theorem~\ref{thm:dirac} and its own lower-bound anchor ($\ce(K_n)\ge n/2$ for even $n$) is a factor of
about $2$; see Section~\ref{sec:open} for the non-regular Dirac problem.
\end{remark}

\section{Maximum degree six: the counterexample and the six-regular case}\label{sec:delta6}

The $n-1$ budget survives $\Delta\le4$ \cite{AABC25} and, as far as anyone can tell, $\Delta\le5$. This
section is about where it dies, and about how much of the $\Delta\le6$ world can be saved.

\subsection{The counterexample}

\begin{theorem}[$\Delta\le6$ with $\ce=n$]\label{thm:counterex}
Let $G$ have vertices $a_1,a_2,a_3,b_1,\dots,b_4$ and edge set consisting of the triangle
$a_1a_2,a_2a_3,a_3a_1$ plus all twelve edges $a_ib_j$; equivalently, $G$ is the union of $K_3$ and
$K_{3,4}$ sharing the vertex set $\{a_1,a_2,a_3\}$. Then $n=7$, $m=15$, $\Delta=6$, and $\ce(G)=7=n$.
\end{theorem}
\begin{proof}[Proof sketch]
An explicit $7$-part certificate --- three cycles of lengths $4,3,4$ and four single edges --- covers each
of the $15$ edges exactly once, so $\ce\le7$. For the lower bound, each $b_j$ has odd degree $3$, so at
least one single edge is incident to each $b_j$ (Lemma~\ref{lem:parity}); any decomposition with at most
$6$ parts therefore has $c+s\le6$ with $s\ge4$, leaving the cases
$(c,s)\in\{(0,4),(0,5),(0,6),(1,4),(1,5),(2,4)\}$. The cases with $c\le1$ cannot cover $15$ edges with
cycles of length at most $7$. For $c=2,s=4$: exactly one single edge is taken at every $b_j$, so the two
cycles must supply all eight remaining edge-ends at the $b$-vertices. The lengths must be $\{4,7\}$ or
$\{5,6\}$. A $7$-cycle visits seven distinct vertices, hence all four $b_j$'s, demanding eight
$b$-side ends on seven edges --- impossible. For $\{5,6\}$: a $6$-cycle is purely bipartite and touches
three $b_j$'s, while the $5$-cycle, which then contains a triangle edge and only three $a$-vertices are
available, needs two \emph{untouched} $b_j$'s --- of which only one remains.
\end{proof}
\proofsrc{Complete proof in campaign records, \texttt{campaign/lane\_05\_delta6.md} Thm.~B. Verified four
independent ways: a DP solver, branch-and-bound, a feasibility enumeration at $k=5,6$, and an independently
written solver with zero shared code (\texttt{campaign/prepub\_independent\_check.py}), plus per-edge
certificate checking (Appendix~A).}

\begin{corollary}\label{cor:d6fail}
``$\Delta\le6\Rightarrow\ce\le n-1$'' is false. Trees show that $n-1$ is attained already at $\Delta\le1$,
and Theorem~\ref{thm:counterex} shows that $n$ is attained at $\Delta=6$; so within $\Delta\le6$ the
$0$-slack behaviour of trees persists at the top degree, and $n$ is the natural next ceiling.
\end{corollary}

\subsection{Eulerian graphs of maximum degree six}

\begin{theorem}[Eulerian graphs of maximum degree $6$]\label{thm:d6even}
Let $G$ be Eulerian with $\Delta(G)\le6$. Then:
(i) $\ce(G)\le 3\lfloor n/3\rfloor\le n$;
(ii) if $n\not\equiv0\pmod3$, or if $G$ is not $6$-regular, then $\ce(G)\le n-1$.
In particular, $\ce(G)\ge n$ forces $n\equiv0\pmod3$ and $G$ $6$-regular.
\end{theorem}
\begin{proof}[Proof sketch]
By \cite{AABC25} (Thm.~1.5: an even graph of maximum degree $\Delta$ decomposes into $\Delta/2$
$2$-regular subgraphs), $E(G)$ splits into $F_1,F_2,F_3$. Each $F_i$ is a union of vertex-disjoint cycles
of length at least $3$, so $F_i$ contributes at most $\lfloor n/3\rfloor$ parts, and the cycles of the three
factors concatenate to a decomposition. If $n=3m'\pm1$ this already gives $\le n-1$. If $G$ is not
$6$-regular, some vertex has degree at most $4$ and is missed by one of the factors, saving one further
part.
\end{proof}
\proofsrc{Complete proof in campaign records, \texttt{campaign/lane\_05\_delta6.md} Thm.~A.}

So everything funnels into one case: $6$-regular, $n\equiv0\pmod3$. Theorem~\ref{thm:d6even} leaves it at
$\ce\le n$; we now close it at $n-2$. Throughout, $6$-regular graphs have $n\ge7$, so under
$n\equiv0\pmod3$ we have $n\ge9$; write $n=3m$.

\begin{lemma}[Petersen's $2$-factorisation, self-contained form]\label{lem:petersen2f}
Every $2k$-regular graph decomposes into $k$ edge-disjoint spanning $2$-factors.
\end{lemma}
\begin{proof}[Proof sketch]
Orient each component along an Euler tour; every vertex gets in-degree $=$ out-degree $=k$. Split each
vertex into an out-copy and an in-copy and view the arcs as edges of a $k$-regular bipartite graph; by
Hall's theorem this graph has a perfect matching, and removing perfect matchings one at a time gives $k$
edge-disjoint perfect matchings (K\"onig). Each matching is a $1$-in/$1$-out spanning digraph without
$2$-cycles (each edge is oriented once), hence a union of directed cycles of length at least $3$; its
underlying undirected graph is a spanning $2$-factor.
\end{proof}
\proofsrc{Complete proof in campaign records, \texttt{campaign/r2\_01\_aabc25\_struct.md} \S1, Lemma 0;
classical, included for completeness.}

\begin{lemma}[Rigidity of triangle-forced $4$-regular graphs]\label{lem:fourreg}
Let $G$ be a $4$-regular simple graph with $|V(G)|\ge5$ in which every $2$-factor is a union of triangles.
Then there are a cubic bipartite graph $B$ and an isomorphism $G\cong L(B)$. In particular $G$ has two
edge-disjoint triangle $2$-factors, and any two of its triangles from different factors share at most one
vertex.
\end{lemma}
\begin{proof}[Proof sketch]
Suppose every $2$-factor of $G$ is a union of triangles, and argue in steps.
(1) By Lemma~\ref{lem:petersen2f} with $k=2$, $G=F\sqcup F'$ with $F$ a $2$-factor; then $F'=G-E(F)$ is
also a $2$-factor, and both are triangle $2$-factors. Write $n=3m$, $m\ge2$.
(2) Crossing pattern. If $a\ne b$ lay in both a triangle $\tau$ of $F$ and a triangle $\sigma$ of $F'$,
then, triangles being cliques in a simple graph, the edge $ab$ would lie in $E(F)\cap E(F')$, which is
empty; so $|V(\tau)\cap V(\sigma)|\le1$. Each vertex lies in exactly one triangle of each factor. If two
vertices of $\tau$ sat in the same $F'$-triangle, that triangle would share two vertices with $\tau$;
hence the three vertices of $\tau$ lie in three \emph{distinct} $F'$-triangles, and symmetrically for
$\sigma$. (When $m=2$ this is already a contradiction, so $m\ge3$.)
(3) Incidence graph. Build a bipartite graph $B$ with sides the triangles of $F$ and of $F'$, and one edge
of $B$ for each vertex $v$ of $G$, joining $\tau(v)$ to $\sigma(v)$. Step (2) says $B$ is simple and
cubic; it is bipartite by construction. Edges of $G$ are exactly the pairs of vertices sharing a triangle,
i.e.\ pairs of $B$-edges sharing a $B$-endpoint, so $G\cong L(B)$.
\end{proof}
\proofsrc{Complete proof in campaign records, \texttt{campaign/r2\_01\_aabc25\_struct.md} \S2, Lemma 1,
steps (1)--(3); full proof in campaign records. The structural reduction appears to have no precedent; the
exhaustive verification below covers every potential counterexample precisely because of it.}

\begin{proposition}[The two complementary $2$-factors of a line graph]\label{prop:antiphi}
Let $B$ be a cubic bipartite simple graph, let $W\cong L(B)$, let $M$ be a perfect matching of $B$, and let
$\sigma$ be an orientation of $M$, with $-\sigma$ the reversed orientation. Then $E(\Phi(\sigma))$ and
$E(\Phi(-\sigma))$ \emph{partition} $E(W)$, and both $\Phi(\sigma)$ and $\Phi(-\sigma)$ are triangle-free
$2$-factors of $W$ (all cycles of length $\ge4$).
\end{proposition}
\begin{proof}[Proof sketch]
For each $z\in V(B)$ let $\{e_z,f_z,g_z\}$ be the three incident edges, $e_z$ being the matched one, and
let $P_z$ be the three pairs of edges sharing $z$. Since $B$ is simple, each edge of $W=L(B)$ is a pair of
$B$-edges sharing a \emph{unique} endpoint, so the $P_z$ partition $E(W)$. Now read pointwise: at a tail of
$\sigma$, $\Phi(\sigma)$ takes the pairs $\{e_zf_z,e_zg_z\}$; at a head, it takes $\{f_zg_z\}$; and
$\Phi(-\sigma)$ does the opposite. In both cases the two factors take distinct pairs summing to three per
$P_z$, so their edge sets partition $E(W)$. A degree check at each $B$-edge --- a matched edge gets $2+0$
neighbours, an unmatched edge $1+1$ --- shows both are $2$-regular and spanning. For triangle-freeness: a
triangle of $L(B)$ corresponds to three $B$-edges meeting at one point, and at that point one of the three
pairs is precisely the pair not taken by the factor in question; $B$ has no triangles of its own, being
bipartite.
\end{proof}
\proofsrc{Complete proof in campaign records, \texttt{campaign/r2\_01\_aabc25\_struct.md} \S2,
Proposition 6 and steps (4)--(8); full proof in campaign records. Positioning: the \emph{existence}
statement --- a cubic graph has a perfect matching iff its line graph admits a triangle-free
$2$-factorisation --- is an observation of Kouider--Sabidussi \cite{KS95}, and the case proved here (cubic
bipartite $B$) is a direct special case of it via Hall's theorem; we claim no priority for it. The explicit
pointwise partition $\Phi(\sigma)\sqcup\Phi(-\sigma)$, its ``for every orientation'' quantifier, and the
counting corollary below are the contribution claimed here. See also Bertram--Hor\'ak \cite{BH97} for a
polynomial decision algorithm for the general $4$-regular case.}

\begin{theorem}[Six-regular graphs, strengthened bound]\label{thm:sixreg}
Every $6$-regular simple graph $G$ with $|V(G)|=n=3m\ge9$ satisfies $\ce(G)\le n-2$.
\end{theorem}
\begin{proof}[Proof sketch]
Take any $2$-factorisation $G=F\sqcup F'\sqcup F''$ (Lemma~\ref{lem:petersen2f}, $k=3$).
Case 1: $G$ has no triangle $2$-factor. Then none of the three factors is a union of triangles, so each
contains a cycle of length $\ge4$ and has at most $1+\lfloor(n-4)/3\rfloor=m-1$ components; the total is at
most $3m-3=n-3$.
Case 2: $G$ has a triangle $2$-factor $S$, with exactly $m$ parts, and $H:=G-E(S)$ is a $4$-regular
spanning simple graph on $n\ge9$ vertices.
--- If $H$ has a triangle $2$-factor $T$, then $S$ and $T$ are edge-disjoint triangle $2$-factors of a
$4$-regular spanning subgraph $W=S\sqcup T$, and the rigidity lemma gives $W\cong L(B)$;
Proposition~\ref{prop:antiphi} splits $W$ into two triangle-free $2$-factors with at most $m-1$ parts each,
while $R:=G-E(W)$ is a $2$-factor with at most $m$ parts. Total: $(m-1)+(m-1)+m=3m-2=n-2$.
--- If $H$ has no triangle $2$-factor, take any $2$-factorisation $H=T\sqcup U$: neither is a triangle
union, so each has at most $m-1$ parts, and the total is $m+(m-1)+(m-1)=n-2$.
\end{proof}
\proofsrc{Complete proof in campaign records, \texttt{campaign/r2\_01\_aabc25\_struct.md} \S3, Theorem
2$'$; full proof in campaign records. Numerics (Appendix~A): exhaustive verification of the structural
lemma's full parameter space for $s\le5$ ($832{,}944$ combinations) and the $L(B)$ family $s=3..6$
($299{,}265$ instances); all $30{,}016$ labelled $6$-regular graphs on $9$ vertices satisfy the bound with
room to spare --- each has exact $\ce=3$.}

\begin{theorem}[Conjecture 6.2 of \cite{AABC25}]\label{thm:sixregconj}
Every $6$-regular simple graph decomposes into three edge-disjoint spanning $2$-factors, one of which has
a component with at least $4$ vertices.
\end{theorem}
\begin{proof}[Proof sketch]
Three cases. If $3\nmid n$, no triangle $2$-factor exists (it would force $3\mid n$), so any
$2$-factorisation from Lemma~\ref{lem:petersen2f} has a non-triangle factor, which contains a cycle of
length $\ge4$. If $3\mid n$ and $G$ has no triangle $2$-factor, the same applies to every factor of any
$2$-factorisation. If $3\mid n$ and $G$ has a triangle $2$-factor $S$, then $|V(G)|=n\ge9$ and
$H:=G-E(S)$ is $4$-regular on at least $5$ vertices; by the contrapositive of
Lemma~\ref{lem:fourreg}, $H$ has a $2$-factor $S'$ containing a cycle of length $\ge4$, and
$S\sqcup S'\sqcup T$ with $T:=H-E(S')$ is the required decomposition.
\end{proof}
\proofsrc{Complete proof in campaign records, \texttt{campaign/r2\_01\_aabc25\_struct.md} \S3, Theorem 3
(the three-case decision written out in full); full proof in campaign records.}

\begin{corollary}[Six-regular graphs]\label{cor:sixregall}
Every $6$-regular simple graph $G$ on $n$ vertices satisfies $\ce(G)\le n-1$. If $n\equiv0\pmod3$ this
holds in the stronger form $\ce\le n-2$ (Theorem~\ref{thm:sixreg}); if $n\equiv2\pmod3$ it holds in the
form $\ce\le n-2$ as well, via $3\lfloor n/3\rfloor=n-2$.
\end{corollary}

\begin{corollary}[Even graphs of maximum degree six]\label{cor:evend6}
Every even graph with $\Delta\le6$ satisfies $\ce(G)\le n-1$: apply Theorem~\ref{thm:d6even}(ii) away from
the exceptional case and Corollary~\ref{cor:sixregall} in it. Together with Theorem~\ref{thm:counterex}
this locates the failure of the $n-1$ budget at $\Delta\le6$ exactly in the presence of odd degrees.
\end{corollary}

\begin{remark}[Tightness is open, and the truth is far lower]\label{rem:sixregtight}
We do not know whether $n-1$ is attained in the even world: the path $P_n$ has $\ce(P_n)=n-1$ but is not
an even graph, so it is no witness here. The computational record suggests the truth is far below the
proof: all $6$-regular graphs on $9$, $10$ vertices have exact $\ce=3$, and $n=12$ has $7{,}849$
isomorphism classes whose decomposition ledgers reach $n-2=10$ pieces in only $5$ classes --- and even
those have exact $\ce=3$. The proof closes at $n-2$ through the ledger, not through the truth. At degree
$8$ the analogous bound is immediate from the all-$2$-factors-isomorphic analysis of \cite{AAFJLS04}, as
noted in \cite{AABC25}; degree $6$ is the last even degree where a structural argument of the kind given
here is needed. Whether the two residual configurations of Theorem~\ref{thm:sixreg}'s case analysis can be
pushed to $n-3$ is registered as an open problem in Section~\ref{sec:open}.
\end{remark}

\subsection{The general $\Delta\le6$ question}

\begin{theorem}[General $\Delta\le6$ graphs]\label{thm:d6tjoin}
If $\Delta(G)\le6$ and $\tau=\tau(G)$, then $\ce(G)\le\tau+\lfloor(m-\tau)/3\rfloor$ with $m\le3n$.
Consequently: $\tau\le 3n/4\Rightarrow\ce\le 3n/2$; $\tau\le n/2\Rightarrow\ce\le 7n/6$; and always
$\ce\le 5n/3-2/3$.
\end{theorem}
\proofsrc{Complete proof in campaign records, \texttt{campaign/lane\_05\_delta6.md} Thm.~D, via
Lemma~\ref{lem:tjoin} with $m\le3n$.}

\begin{remark}[An honest gap]\label{rem:danger}
We do not claim $\ce\le n$ for all $\Delta\le6$. The even case is now closed
(Corollary~\ref{cor:evend6}), but when the graph has odd degrees and $\tau\in(3n/4,\,n-1]$ with
$m>3n-2\tau$ --- long $T$-joins on dense graphs --- no proof of $\ce\le n$ is known to us; this mirrors the
``tree-absorbing rewind'' obstruction identified for $\Delta=5$ in the records. The counterexample of
Theorem~\ref{thm:counterex} sits exactly in such a regime: the extremal accounting cannot close, yet the
true value is $\ce=n$.
\end{remark}

\begin{remark}[The $K_{3,n-3}$ mechanism]\label{rem:k3n}
For $n\ge7$, $\ce(K_{3,n-3})\ge n-1$: the $n-3$ vertices of degree $3$ force $n-3$ single edges, and the
remaining $2(n-3)$ edges exceed the capacity $n$ of a single cycle, forcing at least two cycle parts; the
exact values $\ce(K_{3,4})=6$, $\ce(K_{3,5})=7$, $\ce(K_{3,6})=8$ all equal $n-1$. The case $n=6$ is a
genuine exception: $\ce(K_{3,3})=4=n-2$, because the six residual edges fit in a single Hamilton cycle.
This parity mechanism, ``upgraded'' by the triangle of Theorem~\ref{thm:counterex}, is the principal source
of $\Delta\le6$ lower bounds.
\end{remark}

\section{Odd degrees: the cubic case of Conjecture 6.4}\label{sec:odd}

The paper \cite{AABC25} also studies a relaxation. An \emph{$\fre$-decomposition} of $G$ is a partition of
$E(G)$ into blocks, each block being the edge set of a $2$-regular subgraph (a union of vertex-disjoint
cycles) or a single edge, and $\fre(G)$ is the minimum number of blocks. Every $\ce$-decomposition is an
$\fre$-decomposition, so $\fre(G)\le\ce(G)$; what the relaxation buys is that several cycles may share one
block. Conjecture 6.3 of \cite{AABC25} proposes $\fre(G)=O(n)$ for all graphs, and Conjecture~6.4 there
proposes $\fre(G)\le n-1$ for every $(2k+1)$-regular graph. The even-regular case is settled there
(Thm.~1.5: $\fre(G)=\Delta/2$ for an even graph of maximum degree $\Delta$), and the same paper observes
that a perfect matching settles the odd-regular case: taking the matching as single edges, $\fre\le n/2+k$.
The open core is thus the odd-regular graphs \emph{without} a perfect matching. This section settles the
cubic case completely.

\begin{theorem}[Reduction identity]\label{thm:frered}
For every graph $G$ with $m$ edges,
\[
\fre(G)\;=\;\min_{H}\ \bigl(m-|E(H)|+\Delta(H)/2\bigr),
\]
where the minimum is over even subgraphs $H\subseteq G$.
\end{theorem}
\begin{proof}[Proof sketch]
($\le$) Given $H$, the construction behind \cite{AABC25} Thm.~1.5 produces $\Delta(H)/2$ spanning
subgraphs of $H$, each of every-vertex-degree $0$ or $2$; together with the $m-|E(H)|$ edges of $G-H$ as
single edges this is an $\fre$-decomposition with the displayed count. ($\ge$) From an optimal
$\fre$-decomposition with $t$ two-regular blocks and $s$ single edges, let $H$ be the union of the blocks.
A vertex lying in $r_v$ blocks has $\deg_H(v)=2r_v$, so $H$ is even and $\Delta(H)/2=\max_v r_v\le t$;
hence $s+\Delta(H)/2\le s+t=\fre(G)$.
\end{proof}
\proofsrc{Complete proof in campaign records, \texttt{front\_tre/attack\_tre.md} \S3, Theorem A; full proof
in campaign records.}

Let $p(G)$ denote the minimum size of a parity subgraph with $T=V(G)$, and let $p_2(G)$ denote the maximum
size of an even subgraph. Two immediate consequences of Theorem~\ref{thm:frered}: if $G$ is $(2k+1)$-regular
then $\fre(G)\le p(G)+k$ (remove a minimum parity subgraph); and on cubic graphs the formula sharpens to an
equality, since $G-F$ has degrees in $\{0,2\}$:

\begin{theorem}[Cubic exact formula]\label{thm:cubicexact}
If $G$ is cubic, then $\fre(G)=p(G)+1$, and $p(G)=m-p_2(G)$; consequently Conjecture~6.4 of
\cite{AABC25} for cubic $G$ is equivalent to $p_2(G)\ge n/2+2$.
\end{theorem}
\begin{proof}[Proof sketch]
For a minimum parity subgraph $F$, the complement $G-F$ is an even subgraph of a subcubic graph, hence a
disjoint union of cycles --- a single two-regular block when nonempty --- giving $\fre\le p+1$. Conversely
$G$ contains a cycle $C$, and $G-E(C)$ is an odd factor of size $m-|C|\le m-3$, so $p\le m-3<m$ and the
minimum in Theorem~\ref{thm:frered} is never attained with $H=\emptyset$; the exact value is $p+1$. For
$p=m-p_2$: the complement of a maximum even subgraph is a minimum $T$-join and vice versa.
\end{proof}
\proofsrc{Complete proof in campaign records, \texttt{front\_tre/attack\_tre.md} \S3, Corollary A3; full
proof in campaign records.}

The perfect-matching case of the cubic conjecture is now immediate ($\fre\le n/2+1\le n-1$ for $n\ge4$),
and this is the observation already recorded in \cite{AABC25}. The substance is the graphs without a
perfect matching. Here the Petersen theorem in its contrapositive form is the structural entry point: a
cubic graph without a perfect matching has a bridge. Bridges do not occur inside blocks, and $T$-joins
decompose over the block tree --- which is what the next two results exploit.

\begin{lemma}[Bridge decomposition]\label{lem:bridgedecomp}
Let $G$ be connected cubic with bridge set $Br$, $B=|Br|$, and let $X_1,\dots,X_k$ be the components of
$G-Br$, each bridgeless and subcubic; write $b_v$ for the number of bridges at $v$ and let
$T_i=\{v\in X_i: b_v\ \mathrm{even}\}$. Then $|T_i|$ is even, and
\[
p(G)\;=\;\sum_i q(X_i)\;+\;B,\qquad q(X_i):=\min\{|F|:F\ \text{a $T_i$-join in $X_i$}\}.
\]
\end{lemma}
\begin{proof}[Proof sketch]
Every bridge lies in every odd spanning factor: if a bridge $e$ cuts off $A$, then $|A|$ is odd
(since $3|A|-1=2e(A)$), and $F\cap E(A)$ would have all $|A|$ vertices odd --- impossible by the handshake
lemma. So $B$ edges are forced. Within $X_i$, the requirement $\deg_F\equiv1\pmod2$ becomes
$\deg_{F\cap X_i}\equiv1+b_v\pmod2$, i.e.\ $F\cap X_i$ is a $T_i$-join; and $|T_i|$ is even because
$3n_i-\sum_v b_v=2e(X_i)$ forces $\sum b_v\equiv n_i\pmod2$. Both directions are now immediate, and the
minima add.
\end{proof}
\proofsrc{Complete proof in campaign records, \texttt{front\_tre/attack\_tre.md} \S4, Theorem C; full proof
in campaign records.}

\begin{lemma}[Piece lemma; Choi--Kim--Kostochka--Park--West]\label{lem:piece}
Let $X$ be connected, bridgeless and subcubic with $n\ge3$ vertices. Then $p_2(X)\ge n/2+3/2$. More
generally, in the setting of Lemma~\ref{lem:bridgedecomp}, writing $\Phi(X):=n-q(X)$ and $b$ for the
number of bridge-ends at vertices of $X$, one has $\Phi(X)\ge b/2+3/2$.
\end{lemma}
\begin{proof}[Proof sketch]
On subcubic graphs, maximum even subgraphs are exactly maximum $2$-regular subgraphs counted by edges, so
$p_2$ equals the quantity $f_2$ of \cite{CKKPW19}, whose Theorem 1.3 gives
$f_2\ge n-\max\{0,(d+c-1)/2\}$ in terms of the number $c$ of cut-edges and the $1$-deficit
$d=3n-2m$. Bridgelessness gives $c=0$, and the three branches are: all degrees $2$ (a cycle:
$p_2=n$); all degrees $3$ ($d=0$: $p_2=n$, which is Petersen's $1$-factor theorem); mixed
($d\ge1$ and $m\ge n+1$, giving $p_2\ge n-(d-1)/2=m-n/2+1/2\ge n/2+3/2$). For the corollary form,
$X-H$ is a $T$-join for $H$ maximum even, so $q\le m-p_2$ and
$\Phi=n-q\ge n-m+p_2\ge(3n-2m)/2+3/2=b/2+3/2$, using $\sum_v\deg_X=3n-b=2m$.
\end{proof}
\proofsrc{This is a quoted published theorem plus a conversion derived in campaign records,
\texttt{front\_tre/attack\_tre.md} \S4 (Lemma A = \cite{CKKPW19} Thm.~1.3, whose own proof uses Edmonds'
weighted matching and the theorem of Plesn\'ik; Lemma B = the conversion; full details in campaign
records). The tight example is $K_{2,3}$: $p_2=4=n/2+3/2$.}

\begin{lemma}[Assembly]\label{lem:assembly}
Let $G$ be connected cubic without a perfect matching, with block tree having $k$ blocks ($B=k-1$ bridges):
$r$ non-trivial blocks on $n\ge4$ vertices, $c$ triangle blocks, and $t$ single-vertex blocks. Then
$r\ge c+t+2$, and $\sum_X\Phi(X)\ge B+2$.
\end{lemma}
\begin{proof}[Proof sketch]
Degree conservation over the block tree gives $\sum_X b(X)=2B=2(k-1)$, where $k=r+c+t$ is the number of
blocks; every block has $b\ge1$, and triangle and single-vertex blocks have $b=3$. Hence
$2(r+c+t-1)\ge 3(c+t)+r$, i.e.\ $r\ge c+t+2$. For the ledger: non-trivial blocks give $\Phi\ge b/2+3/2$ by
Lemma~\ref{lem:piece}, triangle blocks give $\Phi=3=b/2+3/2$, single-vertex blocks give $\Phi=1$; summing,
\[
\textstyle\sum_X\Phi\ \ge\ \sum_{\mathrm{pos}}(b_i/2+3/2)+3c+t\ =\ B+1.5r+1.5c-0.5t\ \ge\ B+2,
\]
the last step by $r\ge c+t+2$.
\end{proof}
\proofsrc{Complete proof in campaign records, \texttt{front\_tre/attack\_tre.md} \S4 (assembly lemma); full
proof in campaign records.}

\begin{theorem}[Cubic graphs]\label{thm:cubicpmfree}
Every connected cubic graph without a perfect matching satisfies $p(G)\le n-2$, hence
$\fre(G)=p(G)+1\le n-1$. Together with the perfect-matching case, every cubic graph $G$ on $n\ge4$
vertices satisfies $\fre(G)\le n-1$: Conjecture~6.4 of \cite{AABC25} holds for $k=1$.
\end{theorem}
\begin{proof}[Proof sketch]
Parts never straddle components, so $\fre$ is additive over components and we may assume $G$ connected.
If $G$ has a perfect matching, remove it: $\fre\le n/2+1\le n-1$ for $n\ge4$. Otherwise, by Petersen's
theorem in the contrapositive, $G$ has a bridge, so $k\ge2$. Lemma~\ref{lem:bridgedecomp} writes
$p=\sum_i q_i+B$, and $p=\sum_i q_i+B=n-\sum_X\Phi(X)+B$, so the assembly lemma gives
$p\le n-(B+2)+B=n-2$.
\end{proof}
\proofsrc{Complete proof in campaign records, \texttt{front\_tre/attack\_tre.md} \S4, Theorem D; full proof
in campaign records. Honest scope: the proof relies on the published Theorem 1.3 of \cite{CKKPW19}; the
contribution claimed here is the reduction chain (Theorem~\ref{thm:frered},
Lemma~\ref{lem:bridgedecomp}, Lemma~\ref{lem:assembly}), the exhaustive verification below, and the
structural analysis of the small counterexamples. Computational verification: all $9{,}588$ connected
odd-regular graphs on $n\le12$ vertices (degrees $3$ through $11$), and all connected cubic graphs on
$n=14,16,18$ vertices ($509$, $4{,}060$, $41{,}301$ graphs), with zero violations; all $60{,}077$
bridgeless subcubic graphs on $n\le14$ vertices satisfy Lemma~\ref{lem:piece} with a minimum margin of
exactly $3/2$, attained uniquely by $K_{2,3}$ (Appendix~A).}

\begin{remark}[The smallest imperfect cubic graphs]\label{rem:minpmfree}
The minimum order of a cubic graph without a perfect matching is $16$, and the extremal graph is unique;
the minimality is classical --- it follows from work of Errera and of Chartrand--Goldsmith--Schuster
\cite{CGS79} --- and our contribution is the enumeration confirming uniqueness up to isomorphism for all
orders $\le18$, plus the structure: the extremal graph is three subdivided-$K_4$ blades joined to one
central vertex by three bridges, and it has $p=n/2+1$. In fact every perfect-matching-free cubic graph
found in the exhaustive sweeps has $p=n/2+1$, with the excess over $n/2$ coming entirely from
single-vertex blocks. For $k\ge2$ the whole question reopens: without regularity $3$, ``no perfect
matching'' no longer forces a bridge, and the correct cut structure is that of Tutte sets; a
Tutte-obstruction family of $5$-regular graphs (orders $36$ to $146$) satisfies $p\le n-3$ with room to
spare, and the general case is open.
\end{remark}

\section{Bounded degeneracy}\label{sec:degeneracy}

\begin{theorem}[Degeneracy $2$: the exact frontier]\label{thm:deg2}
Let $G$ be $2$-degenerate with $c=c(G)$ components and $s(G)$ non-tree components. Then
$\ce(G)\le n-c-s(G)\le n-1$, with $\ce(G)=n-1$ if and only if $G$ is a tree. In particular every
$K_4$-minor-free (equivalently, treewidth-$\le2$) graph satisfies $\ce(G)\le n-1$.
\end{theorem}
\begin{proof}[Proof sketch]
Induction on connected $G$ via a vertex $v$ with $\deg(v)\le2$. A bridge at $v$ is charged exactly: it
becomes a forced single edge. If the two $v$-edges' other endpoints $u,w$ are adjacent, both edges are
absorbed into the part containing $uw$ at no cost; if not, dissolve $v$ into the chord $uw$ (solvability
preserves $2$-degeneracy), paying at most one. Components add by Corollary~\ref{cor:excess}'s ledger.
Non-tree components pay their extra $1$; hence equality in $\ce\le n-1$ occurs exactly for connected
forests, i.e.\ trees.
\end{proof}
\proofsrc{Complete proof in campaign records, \texttt{campaign/lane\_08\_degeneracy.md} Thm.~2.1 with
Lemma 2.3 (dissolution lemma); $725$ randomized $2$-degenerate instances machine-checked (Appendix~A).}

\begin{theorem}[Degeneracy $3$: an upper bound]\label{thm:deg3ub}
If $G$ is $3$-degenerate and $n\ge3$, then $\ce(G)\le\lfloor(5n-8)/3\rfloor$; connectedness may be replaced
by the sharper form with $n-c(G)$. More generally, $d$-degenerate graphs satisfy
$\ce(G)\le\lfloor((d+2)n-d(d+1)/2-2)/3\rfloor$ and $\ce(G)\le m$.
\end{theorem}
\begin{proof}[Proof sketch]
A degeneracy ordering gives $m\le dn-d(d+1)/2$; combine with Lemma~\ref{lem:tjoin} and $\tau\le n-c$.
\end{proof}
\proofsrc{Complete proof in campaign records, \texttt{campaign/lane\_08\_degeneracy.md} Thm.~3.3 (T3a),
using the independently re-proven Lemma~\ref{lem:tjoin}.}

\begin{theorem}[Exact values for $K_{3,q}$]\label{thm:k3q}
For every $q\ge2$,
\[
\ce(K_{3,q})=\Bigl\lceil \tfrac{4q}{3}\Bigr\rceil .
\]
Consequently $\sup\{\ce(G)/n: G\ \text{\emph{3-degenerate}}\}\ge 4/3$, attained along $q\equiv0\pmod 3$.
\end{theorem}
\begin{proof}[Proof sketch]
Upper bound: split $K_{3,q}$ into blocks $K_{3,3}$ ($C_6$ plus $3$ single edges, $4$ parts), handling
$q=3k,3k+1,3k+2$ with tail blocks $K_{3,4}$ ($6$ parts) and $K_{3,2}=C_4$ plus $2$ single edges
($3$ parts); each case sums to $\lceil4q/3\rceil$. Lower bound (slot--rate lemma): every cycle $C$ in
$K_{3,q}$ has $|C|-1\le\tfrac53|C\cap B|$, and each degree-$3$ vertex of the large side feeds at most one
cycle; hence $\ce\ge 3q-\lfloor\tfrac53 q\rfloor=\lceil 4q/3\rceil$.
\end{proof}
\proofsrc{Complete proof in campaign records, \texttt{campaign/lane\_08\_degeneracy.md} Thm.~3.5 (T3c); an
arithmetic slip in the $q\equiv2$ line of the first draft was corrected in review ($K_{3,2}$ contributes
$3$ parts, not $1$), the total being unchanged.}

\begin{conjecture}[Matching upper bound for $d=3$]\label{conj:t3b}
Every $3$-degenerate graph satisfies $\ce(G)\le\lfloor(4n-6)/3\rfloor$.
\end{conjecture}
\begin{remark}
Conjecture~\ref{conj:t3b} is equivalent, via the peeling ledger, to a $\tau\le(4n-6-m)/2$-type bound on
$3$-cores; $630$ randomized $3$-degenerate instances with exact $\ce$ and exact $\tau$ showed zero
violations. The state of knowledge here deserves emphasis: the lower bound $4/3$ for $\sup\ce/n$ is
proved, with an exact extremal family (Theorem~\ref{thm:k3q}); the matching upper bound is open, and the
best proved upper bound is $5/3$ (Theorem~\ref{thm:deg3ub}).
\end{remark}

\begin{theorem}[Failure at every $d\ge3$]\label{thm:degex}
For every $d\ge3$ there are $d$-degenerate graphs, on unboundedly many vertices, with $\ce(G)>n-1$:
(i) $d=3$: $\ce(K_{3,7})=10=|V(K_{3,7})|$;
(ii) $d=3$, non-bipartite: for $q\ge14$, $G_q=K_{3,q}+xy$ (one edge inside the large side) satisfies
$\deg(G_q)=3$ and $\ce(G_q)\ge(4q-7)/3>|V(G_q)|-1$;
(iii) $d=5$: $\ce(K_{5,11})\ge16>15$;
(iv) in general, for odd $d=2k+1$ and $q$ large,
$\ce(K_{2k+1,q})\ge\frac{3k+1}{2k+1}q-O(1)>n-1$.
\end{theorem}
\begin{proof}[Proof sketch]
All items use the slot--rate lemma: if every cycle $C$ satisfies $|C|-1\le\rho|C\cap B|$, then any
decomposition has at least $m-\rho\sum_{v\in B}\lfloor\deg(v)/2\rfloor$ parts, since each $v\in B$ lies in
at most $\lfloor\deg(v)/2\rfloor$ cycles. In $K_{3,q}$ every cycle has $|C|-1\le\tfrac53|C\cap B|$ and the
large side supplies exactly $q$ slots; adding $xy$ relaxes the rate inequality only within the same
constant, and $\deg(x)=\deg(y)=4$ adds two slots, whence (ii). For (iv) the constant
$(3k+1)/(2k+1)\to3/2$ matches the global Erd\H{o}s lower bound, so these families are extremal
simultaneously for their degeneracy level and for the whole problem.
\end{proof}
\proofsrc{Complete proof in campaign records, \texttt{campaign/lane\_08\_degeneracy.md} Thms.~4.2--4.3, with
machine certificates: for $K_{3,14}+e$ all $3{,}324$ cycles were enumerated for the rate inequality
($\ce\in[17,21]$), for $K_{5,11}$ all $794{,}530$ cycles ($\ce\ge16$); $\ce(K_{3,7})=10$ by exact solver
(Appendix~A).}

\begin{remark}[The degeneracy frontier]\label{rem:frontier}
Within the degeneracy hierarchy the picture is complete: $n-1$ holds for all $2$-degenerate graphs
(Theorem~\ref{thm:deg2}) and fails by arbitrarily large margins at every $d\ge3$ (Theorem~\ref{thm:degex}).
This kills any goal of the shape ``degeneracy $\le6$ $\Rightarrow$ $n-1$ up to finitely many graphs''.
Note the complementarity with Section~\ref{sec:delta6}: the counterexample there lives in the
$\Delta$-bounded world, Theorem~\ref{thm:degex} in the degeneracy-bounded, $\Delta$-unbounded one; the
common enemy is the odd complete-bipartite kernel.
\end{remark}

\section{Extremal examples and the shape of $f(n)$}\label{sec:extremal}

\begin{theorem}[Complete bipartite lower bounds; restated]\label{thm:kst}
Let $1\le s\le t$. Then: (i) if $s$ is odd, $\ce(K_{s,t})\ge t(3s-1)/(2s)$; (ii) if $s$ is even and $t$ is
odd, $\ce(K_{s,t})\ge(2s+t-1)/2$; (iii) if $s,t$ are even, $\ce(K_{s,t})\ge t/2$. In particular
$f(n)\ge(3/2-o(1))n$. No novelty is claimed: case (i) and the $(3/2-o(1))n$ consequence are known
\cite{E83,BM22}, and are re-derived here by a degree-budget argument; cases (ii)--(iii) are footnote-level
refinements recorded in the archive.
\end{theorem}
\proofsrc{Complete proof in campaign records, \texttt{notes/attack\_bipartite.md} Thm.~1.1 (with small-value
table: $\ce(K_{2,5})=4$, $\ce(K_{3,4})=6$, $\ce(K_{3,7})=10$, $\ce(K_{5,5})=7$, exact).}

\begin{theorem}[A non-bipartite exact family]\label{thm:joinstar}
Let $G_c=K_2\vee K_{1,c-1}$ (two universal vertices joined to a star with $c-1$ leaves), $c\ge3$,
$n=c+2$. Then
\[
\ce(G_c)=c-1+\Bigl\lceil\tfrac{c+2}{3}\Bigr\rceil = n-3+\lceil n/3\rceil\approx\tfrac43 n .
\]
\end{theorem}
\proofsrc{Complete proof in campaign records, \texttt{campaign/lane\_03\_constant.md} Thm.~M2 (two-sided
accounting on the savings identity; exact for all $c$, cross-checked at $c=2,\dots,8$ with multiple
realizations and dual certificates).}

\begin{corollary}\label{cor:fn}
For every $n\ge7$ there is an $n$-vertex graph with $\ce(G)\ge n$; i.e.\ $f(n)\ge n$ for all $n\ge7$ (with
$f(10)\ge11$, $f(11)\ge12$). Two independent routes are known: Theorem~\ref{thm:joinstar} and the
$K_{3,q}$ route (Theorem~\ref{thm:k3q} for $n\ge10$).
\end{corollary}
\proofsrc{Complete proof in campaign records, \texttt{campaign/lane\_03\_constant.md} Cor.~N.1.}

\begin{proposition}[Small $4$-regular graphs]\label{prop:4reg}
Every $2$-connected $4$-regular graph on $n\le8$ vertices satisfies $\ce(G)\le n/2-1$.
\end{proposition}
\proofsrc{Complete proof in campaign records, \texttt{campaign/lane\_04\_delta5\_savings.md} C2: $n=6$ has
a unique isomorphism class; $n=8$ follows from Dirac's theorem ($2\delta=8\ge n$, so the graph is
Hamiltonian, giving $1+\lfloor 8/3\rfloor=3$). The case $n\ge10$ is honestly open: it would require a cycle
of length $\ge n/2+6$ in the residual.}

\begin{remark}[Where does $n-1$ bite?]\label{rem:bite}
Three structurally different mechanisms push $\ce$ to its budget: bridges (trees); parity at density (even
complete graphs, Theorem~\ref{thm:complete}); and concentrated odd degree in near-bipartite kernels
($K_{3,4}$, $K_{3,q}$: Theorems~\ref{thm:k3q} and \ref{thm:counterex}). A structural characterisation of
the graphs with $\ce(G)=n-1$ is a central open problem (Section~\ref{sec:open}); the tight-instance census
for $\Delta\le5$, $n\le7$ (seven $2$-connected isomorphism classes) shows that all small tight instances
fall into the third mechanism or are small dense graphs.
\end{remark}

\section{Anatomy of the $\log^*$ barrier}\label{sec:bm}

We recall the skeleton of \cite{BM22} in the form needed below. A graph $G$ is an
\emph{$(\varepsilon,s)$-expander} if $|N_{G-F}(U)|\ge\varepsilon|U|/\log^2 n$ for every $U$ with
$|U|\le 2n/3$ and every $F$ with $|F|\le s$. BM splits the vertex set into parts that are
$(\varepsilon,s)$-expanders, at a cost of $4sn\log n$ deleted edges with $s=\log^{273}n$, then extracts
cycles from each expander, the remainder after each extraction round having average degree at most
$f(d)=O(\log^{274}d)$; iterating against a fresh tower reference (``tower reset'') gives $O(\log^* n)$
rounds and hence $O(n\log^* n)$ total. All statements below are \emph{within this architecture}: they
concern what any refinement of the parameters, or of the single-step density lemma, can achieve. They are
not unconditional lower bounds against every conceivable algorithm.

\begin{proposition}[Fixed-point lock-in]\label{prop:lockin}
Let $f\colon[d_0,\infty)\to[1,\infty)$ be nondecreasing.
\begin{enumerate}
\item[(a)] If $f$ has a (necessarily superconstant) fixed point $d^*>d_0$ with $f(d^*)\ge d^*$ --- as
happens for every $f(d)=(\log d)^c$, $c\ge1$; for $c=274$ the fixed point of $x=274\ln x$ is
$x^*\approx2095.4$ --- then naive iteration $f^{i}(n)$ never drops below $d^*$: the single-step lemma alone
cannot terminate.
\item[(b)] \emph{(Tower-step condition.)} Suppose there is a monotone $g$ and a constant $C>0$ with
$f(Cg(t))\le Cg(t/2)$ for all large $t$. Then any decomposition process whose only per-step information is
``remainder average degree $\le f(\text{current})$'' and which may reset against the original $n$ has depth
$\le\log^* n+O(1)$. The condition is satisfied by polylogarithmic $f$ (this is exactly BM's bookkeeping),
but not in general: for $f(d)=d/2$ the depth is $\Theta(\log n)$, and no reset changes that.
\end{enumerate}
\end{proposition}
\begin{proof}[Proof sketch]
(a) Monotone iteration converges down to the unique fixed point. (b) Induction:
$d_i\le Cg(\log^{[i]}n)$, each step lowering the tower by one level; this abstracts the induction in
\cite{BM22}, verified against the source.
\end{proof}
\proofsrc{Complete proof in campaign records, \texttt{campaign/lane\_12\_expansion.md} Prop.~12.1. Part
(b)'s statement is the reviewer-narrowed form: for general $f$ the claim is false, and the tower-step
hypothesis is load-bearing.}

\begin{proposition}[The phase transition is a mirage]\label{prop:mirage}
Let $\tilde f(d)=2^{(\ln d)^{\theta}}$, $0<\theta<1$. Then the naive iteration depth of $\tilde f$ from $n$
is $\Theta(\ln\ln\ln n/\ln(1/\theta))$, which exceeds $\log^* n$ for every $n$ with $\log^*n\ge4$; and
$\log^*n=o(\log\log\log n)$. Consequently, a single-step lemma with remainder $2^{(\log d)^\theta}$ would
yield $O(n\log\log\log n)$ --- a bound \emph{weaker} than the existing $O(n\log^* n)$. The
``$2^{(\log d)^\theta}$ rung'' sits \emph{between} the CFS rung ($d^{1-\delta}$, depth
$\Theta(\log\log n)$) and the BM rung (polylog, depth $\log^* n$): it is the next rung down, not up.
\end{proposition}
\begin{proof}[Proof sketch]
Set $x_i=\ln\ln\tilde f^{\,i}(n)$; then $x_{i+1}=\theta x_i+o(1)$. For the comparison, unwind the tower:
$\log^*n=k$ forces $n>2\uparrow\uparrow(k-1)$, so $\ln\ln\ln n\gtrsim 2\uparrow\uparrow(k-4)$. Numerically
at $k=4$: $\ln\ln n\approx10.72$, $\ln\ln\ln n\approx2.37$, depth $\approx22.5>4$. The moral: iteration
depth is a dynamical quantity, and must not be inferred from the pointwise size of the remainder map.
\end{proof}
\proofsrc{Complete proof in campaign records, \texttt{campaign/lane\_12\_expansion.md} Prop.~12.2 (constants
corrected in review: crossover $\ln d\approx5607$ for $\theta=0.9$; fixed point $x^*\approx2095.4$, not
$274\ln274$). An earlier, opposite interpretation of an archive phase-transition claim was withdrawn; see
Acknowledgments.}

\begin{proposition}[Self-improvement is impossible by deletion alone]\label{prop:selfrobust}
\begin{enumerate}
\item[(a)] If $G-F$ is an $(\varepsilon,s)$-expander for some $F\subseteq E(G)$, then $\delta(G)\ge s+1$.
Hence no graph of minimum degree $\le s$ can be turned into an $(\varepsilon,s)$-expander by deleting
edges.
\item[(b)] \emph{(Finite witness, downgraded form.)} The path $P_4$ is an $(\varepsilon_0,0)$-expander for
every $0<\varepsilon_0\le 24/25$ (natural logarithms) and has $\delta(P_4)=1$; this is verified by
exhaustive check of all admissible vertex subsets.
\item[(c)] Combining (a) and (b): there exist minimum-degree-$1$ $(\varepsilon_0,0)$-expanders that admit
no edge-deleted $(\varepsilon',s)$-expander for any $s\ge1$. In BM's architecture, vertex splitting plus
edge deletion is a necessity, not a technical choice.
\end{enumerate}
\end{proposition}
\begin{proof}[Proof sketch]
(a) Delete, at a minimum-degree vertex $v$ of $G-F$, all its incident edges; expansion applied to
$U=\{v\}$ is violated unless $\deg_{G-F}(v)\ge s+1$. (b) For $N=4$, $\log^2N=(\ln4)^2\approx1.92$; the
worst $2$-subset has $|N(U)|=1\ge\varepsilon_0\cdot2/1.92$ iff $\varepsilon_0\le0.961\ge24/25$. (c) is
immediate.
\end{proof}
\proofsrc{Complete proof in campaign records, \texttt{campaign/lane\_12\_expansion.md} Prop.~12.3(a),
12.3(b$'$), 12.3(c). The original part (b) claimed an asymptotic family of sparse expanders; it was
withdrawn after a star-shaped counterexample and an unfixable accounting gap were found, and is replaced
here by the finite statement (b) only.}

\begin{remark}[Optimality within the architecture]\label{rem:bmopt}
Within BM's architecture: the parameter $s$ is monotone-constrained along the whole chain (its value
$s=\log^{273}n$ being forced by the square of the expander-splitting threshold); the deletion price
$4sn\log n$ is exact; the three-random-classes partition is essential (merging attempts die on closed walks
that must avoid a forbidden vertex class); and by
Propositions~\ref{prop:lockin}--\ref{prop:selfrobust} the polylog remainder with tower reset cannot be
beaten from inside. A formal intermediate rung $f(d)=2\uparrow\uparrow(\theta\log^*d)$ would give
$O(n\log\log^* n)$, but no implementation path within the BM machine is known. We record these as
architecture-relative findings.
\end{remark}

\section{Concluding remarks and open problems}\label{sec:open}

The six-regular bottleneck that motivated Section~\ref{sec:delta6} is resolved. What remains is a short
list of problems, each with an identified obstruction.

\begin{problem}[Planar graphs]\label{prob:planar}
Is $\ce(G)\le n-1$ for every planar $G$? Verified by exhaustive enumeration for $n\le6$. The
strengthening ``$\ce\le n-2$ for all planar $G$ containing a cycle'' is false ($K_4$; and $1{,}273$
counterexamples at $n=6$), so any proof must exploit savings beyond triangles.
\end{problem}

\begin{problem}[Torus]\label{prob:torus}
Is $\ce(G)\le n-1$ for every toroidal $G$? The bound, if true, is sharp
($\ce(K_{3,4})=\ce(K_{3,5})=\ce(K_{3,6})=n-1$, Theorem~\ref{thm:k3q}) and would subsume the planar problem.
Eulerian toroidal graphs are settled (Theorem~\ref{thm:eulertorus}).
\end{problem}

\begin{problem}[Maximum degree five; general degree six]\label{prob:delta56}
Does $\Delta\le5\Rightarrow\ce\le n-1$ (Problem 6.1 of \cite{AABC25})? And does $\Delta\le6\Rightarrow
\ce\le n$? The even case of the latter is settled at $n-1$ (Corollary~\ref{cor:evend6}); what remains is
the ``long $T$-join'' danger zone of Remark~\ref{rem:danger}.
\end{problem}

\begin{problem}[Below $n-2$ for six-regular]\label{prob:sixreg}
Can the two residual configurations of Theorem~\ref{thm:sixreg} --- a six-regular graph containing two
edge-disjoint triangle $2$-factors, or one whose $4$-regular remainder has no triangle $2$-factor --- be
pushed to $\ce\le n-3$? Both require bypassing the per-$2$-factor ledger, which pins a triangle factor at
$m$ parts. The known truth is far lower ($\ce=3$ at $n=9$), so the question is about the method, not the
constant.
\end{problem}

\begin{problem}[A triangle-free $2$-factor in every $4$-regular graph]\label{prob:trianglefree}
Does every $4$-regular simple graph have a $2$-factor with \emph{no} triangle at all? The proof of
Proposition~\ref{prop:antiphi} settles the line-graph case; computational reconnaissance on random
$4$-regular graphs ($n=6..14$) and on the $L(B)$ family found no obstruction. The decision algorithm of
Bertram--Hor\'ak \cite{BH97} concerns the stronger two-factor version and gives no counting consequence;
the characterisation results of Hartvigsen \cite{Hartvigsen24} for triangle-free $2$-factors in general
graphs should reposition this problem and was not yet integrated at the time of writing.
\end{problem}

\begin{problem}[The $d=3$ constant]\label{prob:d3const}
Prove or refute Conjecture~\ref{conj:t3b}: is
$\sup\{\ce(G)/n : G \text{ is $3$-degenerate}\} = 4/3$, with extremal family $K_{3,3k}$? Equivalent to a
$\tau$-bound on $3$-cores.
\end{problem}

\begin{problem}[The $4$-core]\label{prob:fourcore}
Is every graph whose $4$-core is small ``$n-1$-good''? Neither a proof nor a counterexample is known; the
$K_{\mathrm{odd},q}$ mechanism needs odd degrees that $4$-core candidates lack, making this the recommended
next target of the degeneracy program.
\end{problem}

\begin{problem}[Linear constants]\label{prob:linear}
Does $f(n)\le\lceil3n/2\rceil+O(1)$? Any explicit-constant $O(n)$ bound would be new. The bipartite version
is governed by Haj\'os' conjecture: an upper bound $\lfloor3n/2\rfloor-1$ for bipartite graphs is
equivalent to the bipartite case of Haj\'os' cycle conjecture (1968) and is open; Haj\'os' conjecture has
been verified for all Eulerian graphs on at most $12$ vertices \cite{HNS17}. Known:
$f(n)\ge(3/2-o(1))n$ \cite{E83,BM22}, $f(n)\ge n$ for $n\ge7$ (Corollary~\ref{cor:fn}).
\end{problem}

\begin{problem}[Beyond Dirac regularity]\label{prob:beyond}
For non-regular graphs with $\delta\ge n/2$, is $\ce\le(1+\varepsilon)n$ (or even $\le n$)? Hamilton
decomposition fails off the regular path; the records locate the obstruction in the critical band.
Similarly, for $2$-connected $4$-regular graphs, does $\ce\le n/2-1$ extend to all $n\ge10$?
(Proposition~\ref{prop:4reg}.)
\end{problem}

\begin{problem}[Odd-regular graphs without perfect matchings, $k\ge2$]\label{prob:oddk}
Conjecture~6.4 of \cite{AABC25} is now settled for $k=1$ (Theorem~\ref{thm:cubicpmfree}). For $k\ge2$ the
bridge route dies --- a bridgeless $(2k+1)$-regular graph without a perfect matching is no longer excluded
by Petersen --- and the correct cut structure is that of Tutte sets; a piece lemma for near-regular blocks
is the missing instrument. Computational support: Tutte-obstruction families of $5$-regular graphs satisfy
$p\le n-3$ with large slack.
\end{problem}

\begin{problem}[Tight instances]\label{prob:tight}
Characterise the graphs with $\ce(G)=n-1$ (Remark~\ref{rem:bite}); even a structural dichotomy for
$2$-connected instances would sharpen the block ledger of Section~\ref{sec:blocks}.
\end{problem}

\begin{problem}[Below $\log^*$]\label{prob:belowlogstar}
Find any machinery producing a decomposition of $o(n\log^* n)$ parts; formally the rung
$O(n\log\log^*n)$ exists (Remark~\ref{rem:bmopt}) but has no known implementation. Equivalently: beat the
tower-reset dynamic of Proposition~\ref{prop:lockin}.
\end{problem}

\appendix
\section{Proof status and computational verification}\label{app:verify}

Every theorem in this paper carries, after its statement or proof, a pointer to its complete proof in the
project records. This appendix summarises the computational verification on which several of those
records rely. Three principles governed the verification: every constructive step of a proof is
implemented and checked by a validator independent of the construction logic (degrees, spanningness,
cycle legality, partition completeness); every exhaustive claim is over a \emph{complete} enumeration, not
a sample; and every checker was subjected to negative controls (deliberately injected errors, all
detected). Scripts are maintained with fixed random seeds; three full re-runs of the six-regular suite
produced byte-identical certificates.

\subsection*{General results}

This subsection covers Sections~\ref{sec:blocks}, \ref{sec:planar}, \ref{sec:degeneracy} and
\ref{sec:extremal}; the underlying records are the lane and notes files of the campaign archive.

\begin{itemize}
\item Exact anchors: $\ce(K_4)=3$, $\ce(K_5)=2$, $\ce(K_6)=5$, $\ce(K_7)=3$, $\ce(K_8)=7$ (independent
exact solvers; the $K_8$ value from a $3.3\times10^7$-step search log). $\ce(K_3\cup K_{3,4})=7$, verified
four independent ways (DP solver; branch-and-bound; feasibility enumeration at $k=5,6$; an independently
written solver with zero shared code), plus per-edge certificate checking.
\item Small planar values: $\ce\le n-1$ verified exhaustively for $n\le6$; the $n-2$ strengthening refuted
by $1{,}273$ counterexamples at $n=6$ (Problem~\ref{prob:planar}).
\item Degeneracy: $725$ randomized $2$-degenerate instances machine-checked for
Theorem~\ref{thm:deg2}; $630$ randomized $3$-degenerate instances with exact $\ce$ and exact $\tau$, zero
violations of Conjecture~\ref{conj:t3b}.
\item Machine certificates for Theorem~\ref{thm:degex}: for $K_{3,14}+e$, all $3{,}324$ cycles enumerated
for the rate inequality, $\ce\in[17,21]$; for $K_{5,11}$, all $794{,}530$ cycles enumerated, $\ce\ge16$;
$\ce(K_{3,7})=10$ by exact solver.
\item Exact small values for Theorem~\ref{thm:kst}: $\ce(K_{2,5})=4$, $\ce(K_{3,4})=6$, $\ce(K_{3,7})=10$,
$\ce(K_{5,5})=7$.
\end{itemize}

\subsection*{Six-regular graphs}

This subsection covers Section~\ref{sec:delta6}; the records are \texttt{campaign/r2\_01\_*} and
\texttt{campaign/harden\_round2/}.

\begin{itemize}
\item Structural family, exhaustive: all marked cubic bipartite $B$ with $s=3..6$ vertices per side ---
$1$, $24$, $2{,}040$, $297{,}200$, totalling $299{,}265$ instances (counts match the theoretical
enumeration exactly) --- checked by both the constructive route and an independent brute-force route; in
addition the $\Phi$-mechanism was force-checked on every one of the $299{,}265$ instances.
\item Full parameter space: for $s\le5$, \emph{all} $B$, \emph{all} perfect matchings and \emph{all}
$2^s$ orientations --- $832{,}944$ combinations ($48+3{,}456+829{,}440$) --- zero counterexamples; for
$s=6$, $300$ sampled $B$. (A two-round total of $1{,}183{,}920$ combinations was accumulated across the
two verification passes; the figures $829{,}440$ and $832{,}944$ refer to the $s=5$ single tier and the
$s\le5$ total respectively.)
\item Anti-$\Phi$ partition (Proposition~\ref{prop:antiphi}): verified on the same $299{,}265$ instances
plus $660$ random large $B$ ($s$ up to $40$).
\item Six-regular, $n=9$: all $30{,}016$ labelled graphs (four isomorphism classes of sizes
$280/5{,}040/4{,}536/20{,}160$, matching the theoretical counts) --- decomposition pieces $\le7=n-2$
throughout, the $n-2$ assertion enforced by an independent implementation; exact $\ce=3$ for all four
classes.
\item Isomorphism-class exhaustive, $n=10/11/12$: $21$, $266$, $7{,}849$ classes; maximum piece counts
$8$, $8$, $10$ against bounds $n-1$, $n-1$, $n-2$ --- zero violations. For $n=10$, exact $\ce$ determined
for every class by two independent engines (branch-and-bound and MILP), all equal to $3$. At $n=12$,
exactly $5$ classes reach the ledger value $n-2=10$ (tight ledger witnesses; their exact $\ce$ is still
$3$).
\item Labelled reconnaissance, $n=10$: $890{,}100$ labelled graphs swept under a $600$-second truncation
(a sampling sweep, honestly not exhaustive; superseded for $n=10$ by the class-exhaustive item above).
\item Random six-regular, $n=15/18/21/24$, $2{,}000$ graphs each ($8{,}000$ total, multiple seeds):
maximum piece counts $11/12/13/13$ against bounds $13/16/19/22$. The sampler (circulant start plus
degree-preserving swap MCMC) yields correlated samples; these sweeps are treated as smoke tests, not as
uniform-random evidence.
\item Adversarial families at $s=7$ ($1{,}500$ instances) and $s=8$ ($600$ instances): partition and
double triangle-freeness checks, all pass.
\item Certificates: $18{,}279$ decomposition certificates re-verified by an independent validator, zero
failures; the validator passed $5/5$ negative controls; three full re-runs produced sha256-identical
outputs.
\item Residual risk boundary, stated plainly: the structural danger zone of
Lemma~\ref{lem:fourreg} is exactly the $L(B)$ family; the exhaustive coverage above reaches $s\le6$
(i.e.\ all potential counterexamples on $n\le18$ vertices), and the case $s\ge7$ ($n\ge21$) rests on the
proof alone.
\end{itemize}

\subsection*{Cubic and odd-regular graphs}

This subsection covers Section~\ref{sec:odd}; the underlying records are the files of
\texttt{front\_tre/}.

\begin{itemize}
\item Exhaustive, connected odd-regular, $n\le12$ (degrees $3$ through $11$): $9{,}588$ graphs, zero
violations of the cubic/odd-regular claims; $f_{re}$ exact values additionally determined for all
instances with $m\le20$ ($116$ graphs).
\item Exhaustive, connected cubic: $n=14$: $509$ graphs; $n=16$: $4{,}060$; $n=18$: $41{,}301$ --- zero
violations of Theorem~\ref{thm:cubicpmfree}.
\item Piece lemma census: all $60{,}077$ connected bridgeless subcubic graphs on $n\le14$ vertices satisfy
Lemma~\ref{lem:piece}, with $\min(2p_2-n)=3$: the margin $n/2+3/2$ is attained, and attained
\emph{only} by $K_{2,3}$ (verified up to isomorphism). For even $n\le14$ the margin is at least $4$.
\item Perfect-matching-free census: minimal order $16$, exactly one graph (structure: three
subdivided-$K_4$ blades on one central vertex, $p=n/2+1$); at $n=18$ exactly four graphs, all with the
same blade structure; all perfect-matching-free cubic graphs found have $p=n/2+1$.
\item $5$-regular Tutte-obstruction families (orders $36$ to $146$, $5$ graphs): $p\le n-3$ throughout.
\item Cross-checks: $p$ computed by two independent paths (brute-force pair enumeration vs.\ weighted
matching), $f_{re}$ by two paths (subset DP vs.\ exact-cover branch-and-bound), $p_2$ cross-checked by
ILP; positive anchors ($K_4$, $K_6$ tight at $n-1$; Petersen $=6$; $C_5=1$; $P_4=3$) all recovered.
\end{itemize}

\section*{Acknowledgments}

This work was carried out by a human-directed \emph{multi-agent AI research system}. Independent
``attack lane'' agents proved the theorems in parallel from a shared briefing and were required to make
every constructive claim independently recomputable (exact solvers, explicit certificates, exhaustive
checks); an adversarial \emph{verifier} agent then re-derived the proofs line by line, re-ran the
verification scripts, and cross-checked quotations against local source copies of the cited papers. The
first verification round issued a triage report: eleven results green-lit, eleven items flagged for repair
(all at the level of argument, wording, or constants; none affected a theorem's conclusion), and one claim
confirmed for withdrawal. A blue-team round implemented all repairs, including the replacement of an
erroneous counting argument in Corollary~\ref{cor:eulertripy}, the narrowing of
Proposition~\ref{prop:lockin}(b) to its tower-step form, the correction of an off-by-one in
Remark~\ref{rem:k3n} and of an arithmetic slip in the proof of Theorem~\ref{thm:k3q}, the correction of
two numerical constants in Propositions~\ref{prop:lockin}--\ref{prop:mirage}, and the withdrawal of an
asymptotic construction in Proposition~\ref{prop:selfrobust}, as well as of an earlier claim that a
phase-transition argument yields $O(n\log\log\log n)$ --- a claim whose direction of comparison was
inverted, as Proposition~\ref{prop:mirage} explains. A second, independent round re-reviewed the proof of
the six-regular theorem (verdict: pass, with six minor findings, all repaired), widened the numerical
coverage to the levels recorded in Appendix~A, and re-examined the literature --- where it surfaced the
observation of Kouider--Sabidussi \cite{KS95} and thereby forced the repositioning of
Proposition~\ref{prop:antiphi} from a claimed-first to an explicit-special-case status, as well as the
correction of the citation record of \cite{AABC25} to its second arXiv version. We believe this
adversarial loop is itself a methodological contribution; the complete error ledger is maintained in the
project repository.

\begin{thebibliography}{29}

\bibitem{AABC25} S. Akbari, J. Aloni, A. Beikmohammadi, and A. Clow, \emph{Tight bounds for cycle--edge
decompositions and covers}, arXiv:2509.01901 (2025); cited at version 2 (September 2025). Theorem 1.3 =
$\Delta\le4$; Theorem 1.5 = even graphs; Problem 6.1 and Conjectures 6.2--6.4 in \S6.

\bibitem{AAFJLS04} M. Abreu, R. E. L. Aldred, M. Funk, B. Jackson, D. Labbate, and J. Sheehan,
\emph{Graphs and digraphs with all $2$-factors isomorphic}, J. Combin. Theory Ser. B \textbf{92} (2004),
395--404.

\bibitem{BHKY62} J. Battle, F. Harary, Y. Kodama, and J. W. T. Youngs, \emph{Additivity of the genus of a
graph}, Bull. Amer. Math. Soc. \textbf{68} (1962), 565--568.

\bibitem{BH97} E. Bertram and P. Hor\'ak, \emph{Decomposing $4$-regular graphs into triangle-free
$2$-factors}, SIAM J. Discrete Math. \textbf{10} (1997), 309--317.

\bibitem{BM08} J. A. Bondy and U. S. R. Murty, \emph{Graph Theory}, Graduate Texts in Mathematics 244,
Springer, New York, 2008.

\bibitem{BM22} M. Buci\'c and R. Montgomery, \emph{Towards the Erd\H{o}s--Gallai cycle decomposition
conjecture}, arXiv:2211.07689 (2022).

\bibitem{CGS79} G. Chartrand, D. L. Goldsmith, and S. Schuster, \emph{A sufficient condition for graphs
with $1$-factors}, Colloq. Math. \textbf{41} (1979), fasc.\ 2.

\bibitem{CFS14} D. Conlon, J. Fox, and B. Sudakov, \emph{Cycle packing}, Random Structures Algorithms
\textbf{45} (2014), 608--626.

\bibitem{CKKPW19} I. Choi, R. Kim, A. V. Kostochka, B. Park, and D. B. West, \emph{Largest $2$-regular
subgraphs in $3$-regular graphs}, Graphs Combin. \textbf{35} (2019), 805--813.

\bibitem{CKLOT16} B. Csaba, D. K\"uhn, A. Lo, D. Osthus, and A. Treglown, \emph{Proof of the
$1$-factorization and Hamilton decomposition conjectures}, Mem. Amer. Math. Soc. \textbf{244} (2016),
no.\ 1152.

\bibitem{Di52} G. A. Dirac, \emph{Some theorems on abstract graphs}, Proc. London Math. Soc. (3)
\textbf{2} (1952), 69--81.

\bibitem{EJ73} J. Edmonds and E. L. Johnson, \emph{Matching, Euler tours and the Chinese postman}, Math.
Programming \textbf{5} (1973), 88--124.

\bibitem{E83} P. Erd\H{o}s, \emph{On some of my conjectures in number theory and combinatorics},
Proceedings of the fourteenth Southeastern conference on combinatorics, graph theory and computing (Boca
Raton, 1983), Congr. Numer. \textbf{39} (1983), 3--19.

\bibitem{EG59} P. Erd\H{o}s and T. Gallai, \emph{On maximal paths and circuits of graphs}, Acta Math.
Acad. Sci. Hungar. \textbf{10} (1959), 337--356.

\bibitem{EGP66} P. Erd\H{o}s, A. W. Goodman, and L. P\'osa, \emph{The representation of a graph by set
intersections}, Canad. J. Math. \textbf{18} (1966), 106--112.

\bibitem{EP184} \emph{Erd\H{o}s Problem \#184}, \url{https://www.erdosproblems.com/184} (accessed
September 2026).

\bibitem{Er71} P. Erd\H{o}s, \emph{Some unsolved problems in graph theory and combinatorial analysis}, in:
Combinatorial Mathematics and its Applications (Proc. Conf., Oxford, 1969), Academic Press, London, 1971,
pp.\ 97--109.

\bibitem{Hartvigsen24} D. Hartvigsen, \emph{Finding triangle-free $2$-factors in general graphs},
preprint, February 2024.

\bibitem{HNS17} I. Heinrich, M. V. Natale, and M. Streicher, \emph{Haj\'os' cycle conjecture for small
graphs}, arXiv:1705.08724 (2017).

\bibitem{KS95} M. Kouider and G. Sabidussi, \emph{Factorizations of $4$-regular graphs and Petersen's
theorem}, J. Combin. Theory Ser. B \textbf{63} (1995), 170--184.

\bibitem{Lucas1883} \'E. Lucas, \emph{R\'ecr\'eations math\'ematiques}, vol.\ 2, Gauthier-Villars, Paris,
1883.

\bibitem{Pe91} J. Petersen, \emph{Die Theorie der regul\"aren Graphs}, Acta Math. \textbf{15} (1891),
193--220.

\bibitem{Py85} L. Pyber, \emph{An Erd\H{o}s--Gallai conjecture}, Combinatorica \textbf{5} (1985), 67--79.

\bibitem{Ri74} G. Ringel, \emph{Map Color Theorem}, Grundlehren Band 209, Springer, Berlin, 1974.

\bibitem{SB77} S. Stahl and L. W. Beineke, \emph{Blocks and the non-orientable genus of graphs}, J. Graph
Theory \textbf{1} (1977), 75--78.

\bibitem{TY97} R. Thomas and X. Yu, \emph{Five-connected toroidal graphs are Hamiltonian}, J. Combin.
Theory Ser. B \textbf{69} (1997), 79--96.

\end{thebibliography}

\end{document}
